\subsection{Measures on a Hilbert space}

Let E be a real Hilbert space, in which the scalar product is denoted \((x|y)\) . There exists an isomorphism j of E onto its dual \(E'\) , characterized by the formula \(\langle x,j(y)\rangle=(x|y)\) for x,y in E (TVS, V, §1, No.~7, Th. 3). We will identify E and \(E'\) by means of j. The Fourier transform of a promeasure \(\mu\) on E is therefore a function \(\mathcal{F}\mu\) on E; when \(\mu\) is a

measure, we have

\[
(\mathcal {F} \mu) (x) = \int_ {\mathrm{E}} e ^ {i (x | y)} d \mu (y) \quad (x \in \mathrm{E}).\tag{64}
\]

THEOREM 3 (Prokhorov--Sazonov) --- Let E be a Hilbert space and \(E_{s}\) the space E equipped with the weakened topology. Let \(\mu\) be a promeasure on E, and \(\Phi\) its Fourier transform. The following conditions are equivalent:

\begin{enumerate}
\def\labelenumi{\alph{enumi})}
\item
  The function \(\Phi\) is continuous on \(\mathbf{E}\) for the Sazonov topology.
\item
  For every \(\varepsilon > 0\) , there exists a nuclear positive quadratic form \(\mathbb{Q}\) on \(\mathbf{E}\) such that \(\Phi(0) - \mathcal{R}\Phi \leqslant \varepsilon + \mathbb{Q}\) .
\item
  The promeasure \(\mu\) is a measure on \(\mathbf{E}_s\) .
\end{enumerate}

\(b) \Rightarrow a)\) : This follows from Prop. 8 of No.~8 (cf.~the inequality (54)).

\(a) \Rightarrow c)\) : This follows from Theorem 2 of No.~10.

\(c) \Rightarrow b)\) : Suppose that \(\mu\) is a measure on \(\mathbf{E}_s\) . Let \(\varepsilon > 0\) . For every integer \(n \geqslant 1\) , the set \(\mathrm{B}_n\) of \(x \in \mathbf{E}\) with norm \(\leqslant n\) is a closed subset of \(\mathbf{E}_s\) , and \(\mathbf{E} = \bigcup_{n \geqslant 1} \mathrm{B}_n\) . Therefore there exists an integer \(n \geqslant 1\) such that

\(\mu (\mathbf{E} - \mathbf{B}_n) <   \frac{\varepsilon}{2}.\) The formula

\[
\mathrm{Q} (x) = \frac {1}{2} \int_ {\mathrm{B} _ {n}} (x | y) ^ {2} d \mu (y)\tag{65}
\]

defines a positive quadratic form \(\mathbf{Q}\) on \(\mathbf{E}\) . Set \(\mathrm{C} = \frac{n^2}{2}\mu (\mathrm{B}_n)\) . If \((e_1,\ldots ,e_p)\) is a finite orthonormal sequence in \(\mathbf{E}\) , then

\[
\sum_ {j = 1} ^ {p} (e _ {j} | y) ^ {2} \leqslant \| y \| ^ {2} \leqslant n ^ {2}
\]

for every \(y \in \mathbf{B}_n\) by Bessel's inequality. It follows by integration that

\[
\sum_ {j = 1} ^ {p} \mathrm{Q} (e _ {j}) = \frac {1}{2} \int_ {\mathrm{B} _ {n}} \sum_ {j = 1} ^ {p} (e _ {j} | y) ^ {2} d \mu (y) \leqslant \frac {n ^ {2}}{2} \mu (\mathrm{B} _ {n}) = \mathrm{C},
\]

therefore Q is nuclear.

Moreover, \(1 - \cos t \leqslant \inf \left(2, \frac{t^2}{2}\right)\) for every real number \(t\) , whence

\[
\begin{array}{r l} \Phi (0) - \mathcal {R} \Phi (x) & = \int_ {\mathrm{E}} \bigl (1 - \cos (x | y) \bigr) d \mu (y) \\ & \leqslant \int_ {\mathrm{B} _ {n}} \frac {1}{2} (x | y) ^ {2} d \mu (y) + \int_ {\mathrm{E} - \mathrm{B} _ {n}} 2 \cdot d \mu (y) \\ & <   \mathrm{Q} (x) + \varepsilon \end{array}
\]

for all \(x \in \mathbf{E}\) . Thus \(b)\) is verified.

Q.E.D.

COROLLARY 1. --- Let \(E_{1}\) and \(E_{2}\) be two Hilbert spaces, u a Hilbert--Schmidt mapping of \(E_{1}\) into \(E_{2}\) , and \(\mu\) a promeasure on \(E_{1}\) . Assume that the Fourier transform \(\Phi\) of \(\mu\) is continuous on \(E_{1}\) . Then the promeasure \(\nu = u(\mu)\) is a measure on \(E_{2}\) equipped with the weak topology.

With the identifications of \(E_{1}\) and \(E_{2}\) with their duals introduced in this No., the Fourier transform of \(\nu\) is equal to \(\Phi\circ u^{*}\) , where \(u^{*}\) is the adjoint of u. Now, \(u^{*}\) is a Hilbert--Schmidt mapping of \(E_{2}\) into \(E_{1}\) (Annex, No.~2), and the quadratic form \(y\mapsto\|u^{*}(y)\|^{2}\) on \(E_{2}\) is therefore nuclear. If \((\mathrm{E}_{2})_{\mathcal{S}}\) denotes \(E_{2}\) equipped with the Sazonov topology, \(u^{*}\) is therefore a continuous linear mapping of \((\mathrm{E}_{2})_{\mathcal{S}}\) into \(E_{1}\) , and \(\Phi_{\nu}=\Phi\circ u^{*}\) is continuous on \((\mathrm{E}_{2})_{\mathcal{S}}\) ; Theorem 3 then shows that \(\nu\) is a measure on the space \(E_{2}\) equipped with the weak topology.

COROLLARY 2. --- Let Q be a nuclear positive quadratic form on the Hilbert space E. The Gaussian promeasure \(\Gamma_{Q}\) on E with variance Q is a measure on \(E_{s}\) .

The Fourier transform \(\Phi\) of \(\Gamma_{Q}\) is equal to \(e^{-\mathbb{Q} / 2}\) . Now, \(e^t \geqslant 1 + t\) for every real number \(t\) , whence \(\Phi(0) - \mathcal{R}\Phi \leqslant \mathbb{Q} / 2\) . The condition \(b)\) of Theorem 3 is therefore verified and \(\Gamma_{Q}\) is a measure on \(\mathbf{E}_s\) .

